% --- 1. INFORMAZIONI GENERALI ---
\section{Informazioni Generali}
\begin{itemize}
	\item \textbf{Luogo/Piattaforma:} {Server Discord\textsubscript{G}}
	\item \textbf{Data:} 2026-04-14
	\item \textbf{Orario di inizio:} 14:30
	\item \textbf{Orario di fine:} 17:10
	\item \textbf{Responsabile\textsubscript{G}:} {Roberto M. Doroftei}
	\item \textbf{Scriba:} {Sofia De Blasi}
\end{itemize}

\vspace{0.5cm}
\textbf{Partecipanti presenti:}
\begin{table}[ht]
	\centering
	\renewcommand{\arraystretch}{1.2}
	\begin{tabularx}{\textwidth}{X l}
		\toprule
		\textbf{Nome e Cognome} & \textbf{Matricola} \\
		\midrule
		Sofia De Blasi & 2111014 \\
		Roberto Mariano Doroftei & 2111031 \\
		Filippo Idiometri & 2035617 \\
		Francesco Marcon & 2110990 \\
		Armando Moda Scarati & 2082864 \\
		Anton Ricardo Rupi & 2054304 \\
		\bottomrule
	\end{tabularx}
\end{table}

\vspace{0.5cm}
\textbf{Partecipanti assenti:}
\begin{table}[ht]
	\centering
	\renewcommand{\arraystretch}{1.2}
	\begin{tabularx}{\textwidth}{X l}
		\toprule
		\textbf{Nome e Cognome} & \textbf{Matricola} \\
		\midrule
		Nessuno & - \\
		\bottomrule
	\end{tabularx}
\end{table}

% --- 2. ORDINE DEL GIORNO ---
\section{Ordine del Giorno}
Gli argomenti previsti per la riunione odierna sono:
\begin{enumerate}
	\item Organizzazione e gestione delle sprint\textsubscript{G} review e retrospective
	\item Messa in pari sulla stesure dell'AdR e del PdP
	\item Pianificazione dello sprint successivo e assegnazione dei task\textsubscript{G} e ruoli
\end{enumerate}

% --- 3. RESOCONTO E DISCUSSIONE ---
\section{Resoconto della Riunione}
\subsection{Organizzazione e gestione delle sprint review\textsubscript{G} e retrospective}
Dal confronto emerso in riunione, il gruppo ha concordato
di svolgere Sprint Review e Sprint Retrospective\textsubscript{G} 
all’interno della riunione del martedì, evitando una 
seconda riunione dedicata il lunedì, ritenuta 
difficilmente sostenibile a livello organizzativo.
Durante la Sprint Review il team verifica 
l’avanzamento rispetto alle attività pianificate, 
analizza eventuali difficoltà e valuta la coerenza 
tra ore previste ed effettive, mentre nella Retrospective 
vengono identificate criticità e azioni correttive.
È stato inoltre ribadito che tutte le PR 
devono essere aperte entro la domenica o, al più tardi,
così da permettere al verificatore\textsubscript{G} di completare 
la revisione\textsubscript{G} prima dell’inizio dello sprint successivo. 
Infine, il team ha stabilito di utilizzare la Review 
anche per allinearsi sul lavoro svolto
per garantire che ogni membro sia aggiornato sullo 
stato delle attività.

\subsection{Messa in pari sulla stesure dell'AdR e del PdP}
Durante la riunione il gruppo ha effettuato un 
allineamento sullo stato di avanzamento 
dell’Analisi dei Requisiti\textsubscript{G} (AdR) e del Piano 
di Progetto (PdP), con l’obiettivo di garantire 
una visione condivisa del lavoro svolto e delle 
attività da proseguire nello sprint successivo. 
È emersa la necessità di chiarire la suddivisione 
delle responsabilità, soprattutto per documenti 
estesi come l’AdR, che richiedono un confronto 
collettivo prima della ripartizione delle sezioni. 
Il team ha inoltre discusso 
l’importanza di mantenere aggiornati i 
documenti in modo coerente, definendo convenzioni 
comuni e assicurando che ogni 
membro comprenda lo stato attuale dei contenuti. 
L’allineamento ha permesso di chiarire dubbi, 
uniformare l’approccio redazionale e stabilire le 
priorità per la prosecuzione coordinata dei due documenti.

\subsection{Pianificazione dello sprint successivo e assegnazione dei task e ruoli}
Lo sprint è stato definito con durata dal 
14 al 21 aprile 2026. I ruoli assegnati per questo sprint 
sono: Responsabile (Roberto M. Doroftei), 
Amministratore\textsubscript{G} (Armando Moda Scarati), Verificatore 
(Anton Ricardo Rupi), Analisti 
(Sofia De Blasi, Francesco Marcon e Filippo Idiometri). 
Dato che attualmente i ruoli di progettista\textsubscript{G} e programmatore\textsubscript{G} 
non sono ancora necessari, si è deciso che 
tre membri siano contribuiscano 
come analisti, più specificatamente 
Filippo Idiometri e Francesco Marcon si occuperanno 
dell'AdR, mentre Sofia De Blasi del PdP. 
Si è stabilito, inoltre, che il Responsabile 
redige il verbale e l'Amministratore gestisce 
la creazione delle issue\textsubscript{G}, gestisce la 
comunicazione con l'azienda referente e controlla, 
prima della fine di ogni sprint, che i tempi 
effettivi delle task nel foglio Excel siano stati 
compilati correttamente (il workflow prevede che la 
compilazione avvenga dopo il merge\textsubscript{G} e la chiusura della 
issue).
 

% --- 4. DECISIONI PRESE ---
\section{Decisioni Prese}

\renewcommand{\arraystretch}{1.3}
\vspace{-0.3cm}
\noindent
\begin{longtable}{c p{12cm}}
	\toprule
	\textbf{Codice} & \textbf{Decisione} \\
	\midrule
	\endfirsthead

	\toprule
	\textbf{Codice} & \textbf{Decisione} \\
	\midrule
	\endhead

	\midrule
	\multicolumn{2}{r}{\textit{Continua nella pagina successiva}} \\
	\midrule
	\endfoot

	\bottomrule
	\endlastfoot

	\textbf{VI-1.1} & Svolgere sprint review e sprint retrospective il martedì prima dello sprint planning\textsubscript{G}. \\[4pt]
	\textbf{VI-1.2} & Definire una definition of done standard con checklist e workflow GitHub unificato. \\[4pt]
	\textbf{VI-1.3} & Sospendere la macroarea 1 dell'ADR in attesa di riscontro aziendale e procedere con la macroarea 2. \\[4pt]
	\textbf{VI-1.4} & Inviare una mail all’azienda per chiarimenti su requisiti minimi, opzionali e funzionalità dell’editor. \\[4pt]
	\textbf{VI-1.5} & Confermare la scadenza domenica sera per le pull request\textsubscript{G} e mantenere aggiornamenti costanti via WhatsApp\textsubscript{G}. \\
\end{longtable}

% --- 5. AZIONI (TODO) ---
\section{Azioni da Svolgere (To-Do)}

\renewcommand{\arraystretch}{1.3}
\begin{longtable}{c c p{4.5cm} p{3cm} l}
    \toprule
    \textbf{Codice} & \textbf{Rif. Decisione} & \textbf{Descrizione Task} & \textbf{Assegnatario} & \textbf{Scadenza} \\
    \midrule
    \endfirsthead

    \toprule
    \textbf{Codice} & \textbf{Rif. Decisione} & \textbf{Descrizione Task} & \textbf{Assegnatario} & \textbf{Scadenza} \\
    \midrule
    \endhead

    \midrule
    \multicolumn{5}{r}{\textit{Continua nella pagina successiva}} \\
    \midrule
    \endfoot

    \bottomrule
    \endlastfoot

    % Issue #116 - Verbale interno
    \ghissue{116} & VI-1.5 & Redigere e completare il verbale interno della riunione del 14/04/2026. & Roberto M. Doroftei & 2026-04-19 \\[4pt]

    % Issue #115 - Mail azienda
    \ghissue{115} & VI-1.4 & Preparare e inviare la mail all’azienda per chiarimenti sui requisiti dell’editor. & Armando Moda Scarati & 2026-04-19 \\[4pt]

    % Issue #113 - Analisi dei Requisiti
    \ghissue{113} & VI-1.3 & Proseguire la stesura dell’AdR lavorando sulla macroarea 2. & Filippo Idiometri / Francesco Marcon & 2026-04-19 \\[4pt]

    % Issue #112 - Piano di Progetto\textsubscript{G}
    \ghissue{112} & VI-1.2 & Aggiornare il PdP con DoD, milestone\textsubscript{G} e struttura sprint. & Sofia De Blasi & 2026-04-19 \\[4pt]

\end{longtable}


% --- 6. PROSSIMA RIUNIONE ---
\section{Prossima Riunione}
\begin{itemize}
	\item \textbf{Data:} 2026-04-21
	\item \textbf{Ora:} 14.30
	\item \textbf{OdG preliminare\textsubscript{G}:} 
	\begin{itemize}
		\item Revisione e verifica dei task del primo e secondo sprint;
		\item Allineamento con gli aggiornamenti del referente aziendale;
		\item avanzamento del Piano di Progetto;
		\item avanzamento dell'Analisi dei Requisiti;
		\item pianificazione del terzo sprint e assegnazione dei ruoli.
	\end{itemize} 
\end{itemize}